\section{Craftax: local goals and reward control ask different questions}
\label{sec:craftax}

Craftax-Classic combines movement, resources, inventory, and survival
\citep{matthews2024craftax}. It is a useful starting point because reaching a nearby visual target
and selecting a rewarding sequence of interactions can require different information. We evaluate
one fitted 512-dimensional grid JEPA in two protocols, both using real simulator outcomes.

\paragraph{Short-horizon goal reaching.}
After seven random warm-up actions, six additional random actions generate a reachable goal image.
The planner minimizes distance between its predicted endpoint latent and the encoded goal. It then
executes the selected six actions \emph{open loop}, without replanning. Success means matching the
goal's final player position; it does not require matching inventory or the entire scene. The planner
reaches 121 of 128 goals, compared with 24 for random actions (Table~\ref{tab:craftax}). Goal generation
and execution share simulator randomness. This is evidence for nearby goal reaching under the stated
protocol, not long-horizon task completion.

\paragraph{Reward-driven planning.}
We separately fit reward and value heads on frozen encodings and evaluate 100-step behavior from
128 starts. The six-step CEM controller replans after every real observation. Its mean return is
$0.47$, below $1.45$ for uniform random actions. A myopic reward controller, which does not use
predicted transitions, also underperforms random actions. The reward head's correlation falls from
$0.92$ on its fitting data to $0.35$ on held-out time rows, and value correlation falls from $0.99$
to $0.27$. Poor performance therefore already appears in reward-based action selection without predicted
transitions; the multi-step result concerns the combined head, dynamics, and planning pipeline.

\begin{table}[t]
\centering
\caption{\textbf{Different tasks, one Craftax checkpoint.} Goal reaching uses 128 nearby goals and
six open-loop actions. Reward control uses 128 starts and 100 steps; $\pm$ is the standard deviation
across starts. These are exploratory single-model results, not the standard Craftax achievement
benchmark.}
\label{tab:craftax}
\inputtable{craftax.tex}
\end{table}

This split motivates the controlled studies that follow. Goal-reaching success shows that a model
can guide some actions; it does not certify reward prediction or policy learning in that model.
Conversely, a weak reward pipeline does not establish that its encoder discarded the relevant
state. Appendix~\ref{app:craftax} records the planner and head-fitting details.
